\documentclass[12pt,a4paper]{report}
\usepackage[utf8]{inputenc}
\usepackage{graphicx}
\usepackage{hyperref}
\usepackage{geometry}
\geometry{a4paper, margin=1in}
\usepackage{listings}
\usepackage{color}

\title{
    \vspace{2cm}
    \textbf{Internship Report} \\
    \vspace{0.5cm}
    \Large{Web-Based IFSF Simulator for Cashier Terminals} \\
    \vspace{1cm}
}
\author{Karim Lahyani}
\date{\today}

\begin{document}

\maketitle

\begin{abstract}
This report details the design, architecture, and implementation of a Web-Based IFSF (International Forecourt Standards Forum) Simulator. The project aims to replicate the functionality of a physical Cashier Terminal (Point of Sale) used in petrol stations, facilitating seamless integration testing with external fuel controllers without requiring expensive physical hardware. The simulator is built using a modern technology stack, comprising an Angular frontend, a Node.js backend, and a PostgreSQL database, communicating via TCP sockets using XML-formatted IFSF standard messages.
\end{abstract}

\tableofcontents

\chapter{Introduction}
\section{Context and Motivation}
In the petrol retail industry, hardware components such as fuel dispensers, payment terminals, and Point of Sale (POS) systems must communicate efficiently and securely. The International Forecourt Standards Forum (IFSF) provides a standardized communication protocol to ensure interoperability between these disparate systems. Developing and testing IFSF-compliant systems traditionally requires expensive physical hardware. This project mitigates that challenge by introducing a fully functional software simulator for a Cashier Terminal.

\section{Project Objectives}
The primary objective of this project is to create a robust, web-based simulator capable of:
\begin{itemize}
    \item Generating complex IFSF XML requests (e.g., Card Payments, Loyalty Awards).
    \item Establishing and maintaining raw TCP socket connections with external fuel controllers.
    \item Providing a user-friendly graphical interface for operators to construct transactions dynamically.
    \item Handling asynchronous hardware prompts (e.g., Cashier Terminal YES/NO confirmations) in real-time.
\end{itemize}

\chapter{System Architecture}
\section{Overview}
The application adopts a classic Client-Server architecture designed to separate the user interface from the heavy network and database processing operations.

\section{Frontend (Client Layer)}
The frontend is built using \textbf{Angular}, providing a Single Page Application (SPA) experience. It serves as the visual POS interface where cashiers input transaction data (e.g., amounts, shopping cart items, loyalty details). The frontend communicates with the backend exclusively via RESTful HTTP APIs.

\section{Backend (Middleware and Network Layer)}
The backend is powered by \textbf{Node.js} and \textbf{Express.js}. It acts as a bridge between the Angular frontend and the external IFSF network. Its responsibilities include:
\begin{itemize}
    \item Exposing REST APIs to receive transaction data from the frontend.
    \item Converting JSON payloads into valid IFSF XML strings.
    \item Managing raw TCP connections to send and receive XML streams.
    \item Handling transaction state and timeouts.
\end{itemize}

\section{Database Layer}
A \textbf{PostgreSQL} database is utilized for persistent storage. It logs all incoming and outgoing messages, transaction requests (\textit{request\_info}), POS states, itemized receipts (\textit{sale\_items}), and final responses. This audit trail is critical for executing complex business logic such as transaction rollbacks and loyalty refunds.

\chapter{Functional Features}
\section{Transaction Generation}
The simulator allows the user to construct various types of transactions, including \textit{CardPayment}, \textit{LoyaltyAward}, and \textit{Payment}. The user can input the System Trace Audit Number (STAN), point of payment IDs, and shift numbers.

\section{Shopping Cart and Amount Management}
Operators can dynamically add up to 35 items to a transaction. The system calculates the total amount, pre-authorization amount, and associates specific product codes and tax codes with each item.

\section{Asynchronous Cashier Prompts}
A critical feature of the IFSF standard is the ability for external hardware to interrupt a transaction and prompt the cashier for input. When the backend receives a \textit{CashierTerminal} XML request, it pauses the TCP pipeline. The frontend, which continuously polls the server, detects the pending prompt and renders a modal dialog asking for a YES/NO confirmation. Upon user interaction, the frontend sends an HTTP POST request to the backend, which subsequently resumes the TCP pipeline and sends the corresponding binary XML response (\verb|<InBoolean>1</InBoolean>|).

\section{Loyalty Award Refunds}
The system implements advanced rollback logic for canceled loyalty transactions. When a \textit{LoyaltyAwardRefund} is triggered, the backend queries the PostgreSQL database for the original transaction using the STAN. It identifies the original prices of the items prior to the discount, generates a new shopping cart without the rebate labels, and recalculates the total amount, ensuring financial consistency.

\chapter{Implementation Details}
\section{TCP Socket Management}
The \texttt{tcpHandler.js} module utilizes the native Node.js \texttt{net} module to establish a TCP client. It manages connection retries, timeouts, and data chunking. Incoming binary streams are decoded from ISO-8859-1 to UTF-8 using the \texttt{iconv-lite} library before being parsed.

\section{XML Parsing and Generation}
The simulator relies heavily on the \texttt{xml2js} library. \texttt{xmlGenerator.js} constructs the massive \textit{CardServiceRequest} XML payloads based on the JSON state received from Angular. Conversely, incoming XML from the TCP socket is parsed into JavaScript objects, allowing the backend to extract the \textit{OverallResult} and \textit{TerminalID}.

\chapter{Conclusion}
The Web-Based IFSF Simulator successfully replicates the complex behavior of a physical Cashier Terminal. By abstracting the intricacies of TCP socket management and XML parsing into a robust Node.js backend, and providing a dynamic Angular frontend, the project offers a highly effective and cost-efficient testing environment for IFSF-compliant forecourt systems.

\end{document}
